\documentclass[10pt,a4paper]{amsart}
\usepackage[margin=19mm]{geometry}
\usepackage{amsmath,amssymb,mathtools}
\usepackage[scr=rsfs,cal=euler]{mathalfa}
\usepackage{xfrac}
\usepackage[unicode,hidelinks,bookmarks=false]{hyperref}
\newcommand{\ie}{i.e.\ }
\newcommand{\cf}{cf.\ }
\newcommand{\resp}{respectively\ }
\newcommand{\Subsection}{\subsection}
\providecommand{\nobreakdash}{\nobreak\hbox{-}}
\setlength{\emergencystretch}{2em}
\multlinegap0pt
\usepackage{bm}
\usepackage{bbm}
\usepackage{dsfont}
\usepackage{xspace}
\usepackage{cancel}
\usepackage{mathtools}
\usepackage{cleveref}
\usepackage{amsthm}
\makeatletter
\@ifundefined{thm}{\newtheorem{thm}{Thm}}{}
\@ifundefined{claim}{\newtheorem{claim}[thm]{Claim}}{}
\@ifundefined{lemma}{\newtheorem{lemma}[thm]{Lemma}}{}
\makeatother
\def\Id{\mathrm{Id}}
\newcommand{\cE}{\mathcal{E}}
\newcommand{\cF}{\mathcal{F}}
\newcommand{\cI}{\mathcal{I}}
\title[Diameter and mixing time of the giant component in the perco]{Diameter and mixing time of the giant component in the percolated hypercube}
\author{Michael Anastos, Sahar Diskin, Lyuben Lichev, Maksim Zhukovskii}
\date{Source: arXiv:2510.13348v4 (CC BY 4.0)}
\begin{document}
\maketitle

\noindent\emph{Source and licence.} Diameter and mixing time of the giant component in the percolated hypercube. Version arXiv:2510.13348v4 (CC BY 4.0), distributed under the Creative Commons Attribution 4.0 International licence, as recorded in the arXiv licence metadata for this exact version and shown on the abstract page https://arxiv.org/abs/2510.13348v4. Full rights notice: licenses/arxiv-2510.13348-CC-BY-4.0.txt.\par\smallskip

\noindent\textit{Editorial scope.} Counted unit: the Lemma whose source label is \texttt{lem:diameter\_for a pair} in the article (source0.tex lines 1349--1351, source label \texttt{lem:diameter\_for a pair}), delivered together with the article's own complete original proof of it (source0.tex lines 1352--1436, one \texttt{proof} environment of 85 lines). No mathematics is written, repaired or re-derived: statement and proof are the article's own text line for line. Required context retained uncounted: lem:bootstrap\_7.1 (lines 1322--1324). Every definition, display and label of those lines is retained unchanged. Omitted from this package and not redistributed: 1--1321, 1325--1348, 1437--1795. Nothing inside a retained excerpt is omitted or altered: every retained body is delivered exactly as the article's lines are written, so no omission receipt of any kind accompanies this package. Citations: the retained excerpts carry no citation command at all, so no bibliography entry of the article is transcribed and no reference is redistributed. Reference adaptation: none was needed and none was made; every reference inside the delivered unit resolves inside it, so no unresolved reference or citation is produced. Printed slips kept literal: no printed word, symbol, hypothesis or number of the counted statement or of its proof was altered. Layout and class: the article's own class is \texttt{amsart} with packages including inputenc, enumitem, showkeys, amsmath, amssymb, amsthm, abstract, xcolor, comment, mathtools, graphicx, caption, subcaption, float; the wrapper is the lane's pinned \texttt{amsart} editorial wrapper at 10pt on A4, which restates the theorem-like environments the retained text needs and adds the article's own macro definitions that the retained text needs (\texttt{\textbackslash Id}, \texttt{\textbackslash cE}, \texttt{\textbackslash cF}, \texttt{\textbackslash cI}), with extra wrapper packages (bm, bbm, dsfont, xspace, cancel, mathtools, cleveref). The delivered numbering is the unit's own. Assets: no retained span contains a figure environment, a graphics-inclusion command, a PDF-page inclusion, a table or a TikZ picture; no image file is copied into this package, the package declares no asset dependency and no image file is redistributed with it. Licence: version arXiv:2510.13348v4 of this article is distributed under the Creative Commons Attribution 4.0 International licence, as recorded in the arXiv licence metadata for this exact version and shown on the abstract page https://arxiv.org/abs/2510.13348v4. A full rights notice is delivered at licenses/arxiv-2510.13348-CC-BY-4.0.txt. Characters: every retained excerpt is pure ASCII, so the delivered engine, XeLaTeX with its default Latin Modern text font, reports no missing character.
\begin{lemma}\label{lem:bootstrap_7.1}
Fix $c>1$ and let $p=p(d)=c/d$. There are constants $K_3=K_3(c)>0$ and $K_4=K_4(c)>0$ such that, for every pair of vertices $u,v\in V(Q^d)$, with probability at least $d^{-K_4}$ there is a path of length at most $K_3d$ between $u$ and $v$ in $Q^d_p$. 
\end{lemma}

\begin{lemma}\label{lem:diameter_for a pair}
Fix $c>1$ and let $p=p(d)=c/d$. There is a constant $K_5=K_5(c)>0$ such that the following holds \textbf{whp}: for every pair $(u,v)\in \cI$, if for each $w\in \{u,v\}$ the number of vertices within distance $(\log d)^5$ from $w$ in $Q_p^d$ is at least $\exp((\log d)^4)$, then there is a path of length at most $K_5d$ between $u$ and $v$ in $Q^d_p$.
\end{lemma}

\begin{proof}
Fix a pair of vertices $(u,v)\in \cI$. For a vertex $w$ and a subset $I\subseteq [d]$, we define $Q(w,I)$ to be the cube of dimension $d-|I|$ containing all vertices which agree with $w$ in each coordinate in $I$ (and where the other coordinates vary). 
In addition, denote by $\mathcal{E}_I=\mathcal{E}_I(u,v)$ the event that, for each $w\in \{u,v\}$, the number of vertices within distance $(\log d)^5$ from $w$ in $Q(w,I)_p$ is at least $\exp((\log d)^4)/\max\{|I|,1\}$. Observe that $\cE_{\varnothing}$ is the event that for each $w\in \{u,v\}$, the number of vertices within distance $(\log d)^5$ from $w$ in $Q^d_p$ is at least $\exp((\log d)^4)$.

Set $s=d/(\log d)^2$ and consider a family of disjoint subsets $I_1,I_2,\ldots,I_s$ of $\Id(u,v)$, each of size $d/(3s)=(\log d)^2/3$.
 
\begin{claim}\label{claim:smallerevents}
$\cE_\varnothing\subseteq \bigcup_{j=1}^s \cE_{I_j}$.
\end{claim}
\begin{proof}
Assume that the event $\mathcal{E}_\varnothing$ holds. 
Then, for each $w\in \{u,v\}$, there is a tree rooted at $w$ in $Q^d_p$, denoted by $T_w$, which has depth at most $(\log d)^5$ and contains $\exp((\log d)^4)$ vertices.
For each vertex $x\in V(T_w)$, we denote by $P_x$ the unique path from $x$ to the root $w$ in $T_w$. 
We also denote by $L_w(x)$ the set of coordinates $i$ such that some vertex $x'$ on $P_x$ satisfies $x'(i)\neq w(i)$.
Note that $|L_w(x)|\le e(P_x)\le (\log d)^5$. Furthermore, if $L_w(x)$ does not intersect $I_j$, then $P_x\subseteq Q(w,I_j)_p$.
In particular, for every $x\in V(T_w)$, $P_x\not\subseteq Q(w,I_j)_p$ for at most $(\log d)^5$ indices $j\in [s]$.

For each $w\in \{u,v\}$, denote by $B_w$ the set of indices $j\in [s]$ for which the event $\cE_{I_j}$ fails for $w$. The above implies that, for each $w\in \{u,v\}$,
\[ (\log d)^5 \cdot |V(T_w)|  \ge \sum_{j\in B_w} \left(1-\frac{1}{|I_j|}\right)|V(T_w)|\ge \frac{|B_w||V(T_w)|}{2},\]
where the first inequality is obtained by counting the number of pairs $(I_j,x)$ such that $P_x\not\subseteq Q(w,I_j)_p$ in two different ways: for the upper bound in the left hand side, we fix $x$ and bound the number of choices of $I_j$ by $(\log d)^5$, and for the lower bound in the right hand side, we fix $j\in B_w$ and use the definition of $B_w$ to bound from below the number of vertices in $T_w$ which are not reachable in $Q(w,I_j)_p$ by $|V(T_w)|-|V(T_w)|/|I_j|$.
Hence, $|B_w|\leq 2(\log d)^5$ for each $w\in \{u,v\}$ and therefore, the event $\mathcal{E}_{I_j}$ holds for at least $s-4(\log d)^5\ge 1$ indices $j\in [s]$.    
\end{proof}

Next, we fix $j\in [s]$ and let $L_j\coloneqq L_j(u,v)\subseteq I_j$ be an arbitrary subset of size $|I_j|/2=d/6s=(\log d)^2/6$.
We further define $r=r(d):= \binom{d/6s}{d/12s}\ge 2^{d/12s}$ and denote by $L_{1,j},L_{2,j},...,L_{r,j}$ the subsets of $L_j$ of size $|L_j|/2=d/12s$. 
Also, for each $t\in [r]$, denote by $Q_{t,j}$ the subcube of dimension $d-|I_j|$ with vertex set 
\[\{z\in V(Q^d): z(i)=u(i) \text{ for every }i\in I_j\setminus L_{t,j} \text{ and } z(i)\neq u(i) \text{ for every } i\in L_{t,j}\},\]
noting that the coordinates outside $I_j$ vary and that for every $i\in I_j$, $u(i)=v(i)$. Further observe that the cubes $Q_{1,j},\ldots,Q_{r,j}$ are pairwise vertex-disjoint.

\begin{claim}\label{claim:join_vertices_tocubes}
For every $w\in \{u,v\}$ and $t\in [r]$, the probability that $\mathcal{E}_{I_j}$ occurs and there does not exists a path $P(w,t,j)$ of length at most $d$ from $w$ to $Q_{t,j}$ in $Q^d_p$ which is edge-disjoint from all of the hypercubes $Q_{1,j},Q_{2,j},\ldots Q_{r,j}$ is $o(n^{-4})$.
\end{claim}
\begin{proof}
Fix $w\in \{u,v\}$ and $t\in [r]$.
Reveal the edges in $Q(w,I_j)_p$ and assume that the set of vertices $S_w$ within distance $(\log d)^5$ from $w$ in $Q(w,I_j)_p$ has size at least $\exp((\log d)^4)/|I_j|$ (an event which is implied by $\mathcal{E}_{I_j}$). 
For $z\in S_w$, we denote by $P_{t,z}$ a shortest path from $z$ to $Q_{t,j}$ in $Q^d$. 
By construction, these $|S_w|$ paths are of length $|L_{t,j}|=d/12s$, and edge-disjoint from each of the cubes $(Q_{l,j})_{l\neq t}$. Furthermore, for all distinct $z,z'\in S_w$, if we let $i\in [d]$ be such that $z(i)\neq z'(i)$, then $i\notin \Id(z,z')\supset I_j$. It is clear that all vertices in $P_{t,z}$ agree on coordinates outside of $I_j$. Thus, every vertex on $P_{t,z}$ differs from every vertex 
on $P_{t,z'}$ on their $i$-th coordinate (at least), that is, the $|S_w|$ paths are pairwise vertex-disjoint. Hence, since all paths $P_{t,z}$, $z\in S_w$, are edge-disjoint from $Q(w,I_j)$, and thus their appearance in $Q_p^d$ is independent of $Q(w,I_j)_p$, the probability that none of these $|S_w|$ vertex-disjoint paths belongs to $Q_p^d$ is at most
\begin{equation*}
    (1-p^{d/12s})^{\exp((\log d)^4/6)}\leq \exp\left(-(c/d)^{(\log d)^2/12}\exp((\log d)^4/6)\right) = o(n^{-4}).\qedhere
\end{equation*}
\end{proof}

Now, let us finish the proof of Lemma~\ref{lem:diameter_for a pair}. For $(u,v)\in \cI$, $j\in [s]$, $t\in [r]$, and a constant $K'>0$, denote by $\cF(u,v,j,t,K')$ the event that there exists $u',v'\in Q_{t,j}$ such that each of the following holds:
\begin{itemize}
    \item there are paths of length at most $d$ from $u$ to $u'$ and from $v$ to $v'$ in $Q^d_p$ which do not intersect any of the hypercubes $Q_{1,j},Q_{2,j},\ldots Q_{r,j}$, and
    \item there is a path of length at most $K'd$ between $u'$ and $v'$ in $(Q_{t,j})_p$.
\end{itemize}
Observe that $\cF(u,v,j,t,K')$ implies that there is a path between $u$ and $v$ in $Q^d_p$ of length at most $(K'+2)d$. Thus, for every constant $K'>0$,
the statement of the lemma fails with probability at most 
\begin{align}\label{eq:001}
    \sum_{(u,v)\in \cI}  \mathbb{P}\bigg(\cE_{\varnothing}(u,v)\cap \bigcap_{j\in [s]} \bigcap_{t\in[r]}\cF(u,v,j,t,K')^c\bigg)
\end{align}
Thereafter, by \Cref{claim:smallerevents}, we have that
\begin{align*}
\cE_{\varnothing}(u,v)\cap \bigcap_{j\in [s]} \bigcap_{t\in[r]}\cF(u,v,j,t,K')^c
&\subseteq \bigg(\bigcup_{j'\in [s]}\cE_{I_{j'}}(u,v)\bigg)\cap \bigcap_{j\in [s]} \bigcap_{t\in[r]}\cF(u,v,j,t,K')^c\\
&\subseteq \bigcup_{j'\in [s]}\bigg(\cE_{I_{j'}}(u,v)\cap  \bigcap_{t\in[r]}\cF(u,v,j',t,K')^c  \bigg).
\end{align*}
Hence, the expression in \eqref{eq:001} is bounded above by 
\begin{align}\label{eq:002}
    \sum_{(u,v)\in \cI}  \sum_{j\in[s]}\mathbb{P}\left(\cE_{I_j}\cap \bigcap_{t\in[r]}\cF(u,v,j,t,K')^c\right).
\end{align}

To that end, we estimate the probability of $\cE_{I_j}\cap \cF(u,v,j,t,K')^c$. 
We reveal all the edges of $Q_p^d$ lying outside the cubes $Q_{1,j},Q_{2,j},\ldots Q_{r,j}$. 
By \Cref{claim:join_vertices_tocubes}, for all $t\in [r]$,
with probability $1-o(n^{-4})$, the event $\cE_{I_j}$ holds and there exist paths from $u$ to $Q_{t,j}$ and from $v$ to $Q_{t,j}$ of lengths at most $d$ spanned by the set of revealed edges.
For every $t\in [r]$, denote by $u_t$ and $v_t$ the endpoints of these two paths in $Q_{t,j}$.
Then, \Cref{lem:bootstrap_7.1} implies that there exist constants $K_3',K_4'$ such that we can further connect $u_t$ and $v_t$ in $(Q_{t,j})_p$ by a path of length at most $K_3'd$ with probability at least $d^{-K_4'}$; 
note that these events are independent for different indices $t\in [r]$ since the cubes $Q_{1,j},Q_{2,j},\ldots Q_{r,j}$ are vertex-disjoint. 
Hence,  
\begin{align*}
\mathbb{P}\left(\cE_{I_j}\cap \bigcap_{t\in[r]}\cF(u,v,j,t,K')^c\right)&\le o(rn^{-4})+\left(1-d^{-K_4'}\right)^{r}\\
&\leq  o(n^{-3})+\exp(-d^{-K_4'}r)\\&\le o(n^{-3})+\exp\left(-d^{-K_4'}2^{(\log d)^2/12}\right)=o(n^{-3}),
\end{align*}
where the error term is independent of $u,v,j$ and $t$.

Therefore, by~\eqref{eq:002}, the statement of the lemma fails with probability at most
\[\begin{split}
&\sum_{(u,v)\in \cI} \sum_{j\in [s]} \mathbb P\bigg(\cE_{I_j}(u,v)\cap \bigcap_{t\in [r]} \cF(u,v,j,t,K_3')^c\bigg)\\
&\le \sum_{(u,v)\in \cI} \sum_{j\in [s]} \mathbb P\bigg( \bigcap_{t\in [r]}\cE_{I_j}(u,v)\cap \cF(u,v,j,t,K_3')^c\bigg)=o(|\cI|sn^{-3})=o(1),
\end{split}\]
as required.
\end{proof}

\end{document}
